% 《大学物理简明教程》第5版，解答题4.3.18—4.3.23。
% 原书及PDF第105—106页，均依据扫描页实际图形及参数核对。
\subsection{解答题：波形追踪、干涉与相干条件}

\sourceproblem{4.3.18}{105--106}{105--106}
\problemtext 已知平面简谐波的波函数
$y=A\cos[\pi(4t+2x)]$（SI）。
(1) 写出$t=4.2\,\mathrm s$时各波峰位置、离原点最近的波峰与波峰下次经过原点的时间；
(2) 给出$t=4.2\,\mathrm s$时的波形。
\workings
$\omega=4\pi\,\mathrm{rad/s}$、$k=2\pi\,\mathrm{rad/m}$，
$\lambda=1\,\mathrm m$、$T=0.5\,\mathrm s$，
相位为$\omega t+kx$，故波沿$-x$传播，$u=2\,\mathrm{m/s}$。
波峰条件
\[
\pi(4t+2x)=2n\pi
\quad\Longrightarrow\quad
x_n(t)=n-2t,\quad n\in\mathbb Z.
\]
在$t=4.2$时
\[
\boxed{x_n=n-8.4\ \mathrm m,\quad n\in\mathbb Z}.
\]
离原点最近的波峰为$n=8$对应$x=-0.4\,\mathrm m$，
故$\boxed{|x|_{\min}=0.4\,\mathrm m}$。
波峰经过原点需$4t=2n$，即$t=n/2$；
$t>4.2$的最早值是$\boxed{t=4.5\,\mathrm s}$。
固定$t=4.2$，空间波形可写为
\[
\boxed{y(x,4.2)=A\cos(0.8\pi+2\pi x)}.
\]
\crosscheck 波峰间隔1 m，随时间向左以2 m/s平移；其轨迹$dx_n/dt=-2$。

\sourceproblem{4.3.19}{106}{106}
\problemtext 如图4.3.19，波沿直线向右以$u=20\,\mathrm{m/s}$传播，四点从左到右为$C,B,A,D$，距离$CB=8$ m、$BA=5$ m、$AD=9$ m。已知$A$的振动方程$y_A=3\cos(4\pi t)$。分别以$A$、$B$为坐标原点写出波函数及$C,D$处振动方程。
\workings
$\omega=4\pi$，$k=\omega/u=\pi/5\,\mathrm{rad/m}$。
以$A$为原点、向右为正，
\[
\boxed{y_A(x,t)=3\cos\left(4\pi t-\frac\pi5x\right)}.
\]
$C$位于$x_C=-13$ m、$D$位于$x_D=9$ m，故
\[
\boxed{y_C(t)=3\cos(4\pi t+13\pi/5)},\quad
\boxed{y_D(t)=3\cos(4\pi t-9\pi/5)}.
\]
以$B$为原点，$A$位于$x'=5$ m，故
\[
\boxed{y_B(x',t)=3\cos\left(4\pi t-\frac\pi5(x'-5)\right)}
\]
（这里$y_B$表示\emph{以B为原点的整条波函数}而不专指B点）。
$C$为$x'=-8$ m、$D$为$x'=14$ m，代入重新得到上面的$y_C,y_D$。
\crosscheck 在以$A$为原点的波函数代入$x=0$恢复已给定$y_A$；坐标平移不改变任何实际质点的振动方程。振幅3的单位沿用原书未标注单位，不擅自补成3 m。

\sourceproblem{4.3.20}{106}{106}
\problemtext 两等振幅相干波源$S_1,S_2$，$S_1$至观测点$P$为$30$ cm，$S_1S_2=40$ cm且$\angle S_1=90^\circ$。在此时$S_1$位于波峰、$S_2$位于波谷，传播速率$u=100\,\mathrm{m/s}$，求在$P$发生相长干涉所需的最小频率。
\workings
由直角三角形$S_2P=50$ cm，故两路程差
$\Delta r=0.50-0.30=0.20$ m。
源在同一时刻波峰与波谷，初相差$\pi$；到达点相长条件：
\[
\pi+\frac{2\pi\Delta r}{\lambda}=2n\pi
\quad\Longrightarrow\quad
\Delta r=\left(n-\frac12\right)\lambda.
\]
为了获得\emph{最小正}频率，取满足条件的最大波长
$\lambda_{\max}=2\Delta r=0.40\,\mathrm m$，于是
\[
\boxed{\nu_{\min}=\frac{u}{\lambda_{\max}}
=\frac{100}{0.40}=250\,\mathrm{Hz}}.
\]
\crosscheck 反相波源在零程差点必相消，需额外半波长差抵消反相才能加强。

\sourceproblem{4.3.21}{106}{106}
\problemtext 两相干波源$S_1,S_2$在同一条$x$轴上，间距$d=30\,\mathrm m$；$S_1$为$x=0$，$S_2$为$x=30$ m。两波振幅相等，在$x_1=9$ m及$x_2=12$ m处为\emph{相邻}的干涉静止点，波在介质中向外传播时波长不变。求波长及两源初相差。
\workings
在源间$0<x<d$，两路程为$r_1=x,r_2=d-x$，
故路程差$\Delta r=d-2x$。沿$x$每增加$\Delta x$，干涉相位差变化$2k\Delta x$。
相邻相消位置的相位差相隔$2\pi$，因而
\[
2k(x_2-x_1)=2\pi
\quad\Longrightarrow\quad
\boxed{\lambda=2(x_2-x_1)=6\,\mathrm m}.
\]
在$x_1=9$ m，$r_2-r_1=30-18=12\,\mathrm m=2\lambda$，传播相位差为$4\pi$，模$2\pi$等于零。
为了在此处相消，两源的初相差必须为奇数倍$\pi$：
\[
\boxed{\varphi_2-\varphi_1=\pi\pmod{2\pi}}.
\]
\crosscheck 在$x_2=12$ m处路程差$6$ m$=\lambda$，同样与源反相合成相消；区间内$x=10.5$ m则由相位差额外变化$\pi$而成为相长点。

\sourceproblem{4.3.22}{106}{106}
\problemtext 图4.3.22中的等振幅相干波源$S_1,S_2$位于一条直线，间距为$\lambda/4$；$S_1$的初相位比$S_2$超前$\pi/2$，各自振幅为$A_1$。求$S_1$左侧和$S_2$右侧各点的合振幅及波强。
\workings
取由$S_1$指向$S_2$为$x$正向，$d=\lambda/4$。
对于左侧$x<0$，$r_2-r_1=d$。两波到达该点的相位差
\[
\theta_1-\theta_2
=(\varphi_1-\varphi_2)-k(r_1-r_2)
=\frac\pi2+\frac\pi2=\pi.
\]
故合振幅
$\boxed{A_{\rm left}=0}$，平均波强
$\boxed{I_{\rm left}=0}$（理想相干、等振幅）。
对于右侧$x>d$，$r_1-r_2=d$，于是
$\theta_1-\theta_2=\pi/2-\pi/2=0$，
\[
\boxed{A_{\rm right}=2A_1},\qquad
\boxed{I_{\rm right}=4I_1},
\]
其中$I_1$是单独一列波在该处传播介质中的平均波强。
\crosscheck 平均波强正比振幅平方；此结论假设在直线上两波传播到这些点的振幅相等、无额外几何衰减。

\sourceproblem{4.3.23}{106}{106}
\problemtext 两相干波源$B,C$到观测点$P$的距离分别$BP=0.40$ m、$CP=0.50$ m，传播速率$u=0.20\,\mathrm{m/s}$。
两源振动分别为$y_B=2\times10^{-3}\cos(2\pi t)$与
$y_C=2\times10^{-3}\cos(2\pi t+\pi)$（SI）。
求两波到达$P$的相位差及同时合振动的振幅。
\workings
两源频率均$1\,\mathrm{Hz}$，故波长$\lambda=u/\nu=0.20$ m，$k=2\pi/\lambda=10\pi\,\mathrm{rad/m}$。
在$P$处第二列波相对第一列波的相位差
\[
\Delta\varphi=(\pi-k CP)-(0-k BP)
=\pi-k(CP-BP)
=\pi-10\pi(0.10)=0.
\]
所以
\[
\boxed{\Delta\varphi=0\pmod{2\pi}},\qquad
\boxed{A_P=2A=4\times10^{-3}\,\mathrm m}.
\]
\crosscheck 两源虽然反相，但较长传播路径产生额外$\pi$相移，使得$P$点两波恰好同相；只有求到达点相位才能判断加强/相消。
